\chapter{Conclusions}
\label{cha:conclusioni}
Both classifiers achieved high and comparable performance. The SVC, however, has an advantage in terms of interpretability, as it allows the decision boundary to be expressed through an explicit quadratic equation. Based on the results, the decision boundaries of the two models appear to closely reproduce the separation that emerges from the distribution of outcomes presented in Section~\ref{subsec:relazione-parametri-esiti}. Classification errors are concentrated mainly in the transition region between the two regions, where clusters characterized by the same orbital parameters can reach different outcomes depending on their initial pseudo-random realization. This variability seems to suggest that the semi-major axis and eccentricity, considered as the only input variables, are not sufficient to fully describe the evolution of star clusters.\\ To investigate this aspect further, it may be useful to increase the number of realizations for each orbital configuration and include additional diagnostic quantities among the input variables.\\
\newline A further aspect that emerged during the study concerns cluster dissolution due to their internal dynamical evolution, which can occur even independently of the galactic potential. To limit its influence on the results, the final simulation time was set to $T_{\mathrm{end}}=700$. A possible direction for future work would be to increase the number of stars in each cluster, thereby reducing the effects associated with the finite number of particles and obtaining a more stable description of its internal evolution.\\
\newline A positive result, on the other hand, concerns the Plummer model used to generate the initial cluster conditions. Although its theoretical description is formulated in the limit $N\to\infty$, it proved effective even for clusters containing only 200 stars. This result was supported by the rescaling of the initial velocities, which made it possible to compensate for deviations from virial equilibrium due to sampling with a finite number of stars.\\
\newline Finally, the range of semi-major axes explored represents a further limitation of the \emph{dataset}.\\ It would be interesting to extend this range and verify whether the models can correctly predict the outcomes of configurations characterized by orbital parameters outside the intervals used during training. This analysis would make it possible to evaluate the classifiers' ability to extrapolate the learned relationships and apply them to configurations that have never been simulated before.%
